% MA 214 Discussion 14 -- covers Lecture 21, and doubles as Quiz 4 / Module 4 review
% and the course-wide final review.
%
% Student handout. Page 1 is the Module 4 concept review; the question set runs from page 2.
% Questions are tagged with the Module 4 objective they test; Question 5 spans the whole course.
%
% Compile from inside this folder:  pdflatex Discussion14.tex

\documentclass[11pt]{article}
\usepackage[margin=1in]{geometry}
\usepackage{parskip}
\usepackage{fancyhdr}
\usepackage{array}
\usepackage{booktabs}
\usepackage{enumitem}
\usepackage{amsmath}
\usepackage{tabularx}
\usepackage{listings}

\pagestyle{fancy}
\fancyhf{}
\cfoot{\thepage}
\rhead{Discussion 14}
\lhead{MA 214: Applied Statistics}
\renewcommand{\headrulewidth}{.4mm}

\newcommand{\blankline}[1]{\underline{\hspace{#1}}}
\newcolumntype{L}{>{\raggedright\arraybackslash}X}

\lstset{
  basicstyle=\ttfamily\small,
  columns=fullflexible,
  frame=single,
  framerule=0.4pt,
  breaklines=true,
  showstringspaces=false,
  aboveskip=8pt,
  belowskip=8pt
}

% Section header: bold title over a hairline rule.
\newcommand{\parttitle}[1]{%
  \par\medskip\noindent
  {\large\bfseries #1}\par
  \vspace{-.5em}\noindent\rule{\linewidth}{0.4pt}\par\smallskip}

% Writing space.
\newcommand{\answerspace}[1]{\vspace{#1}}

% Objective tag printed after a question title.
\newcommand{\lo}[1]{{\normalfont\small (#1)}}

% Flush-left question list: the label runs in at the margin and the body wraps
% back to the margin, so nothing is pushed far to the right.
\newlist{questions}{enumerate}{1}
\setlist[questions]{label=\textbf{Question \arabic*.}, wide=0pt, itemsep=1.4em, topsep=.8em}

\newlist{parts}{enumerate}{1}
\setlist[parts]{label=\textbf{(\Alph*)}, leftmargin=1.6em, labelsep=.45em,
  itemsep=.7em, topsep=.5em}

\begin{document}

\noindent\textbf{Name:}\blankline{2.3in}\hfill\textbf{BUID:}\blankline{1.6in}

\begin{center}
  {\Large\bfseries Discussion 14: Bayesian Regression, and the Whole Course}
\end{center}

\noindent\textbf{Goals:} run Module 4's last model, then the module, then the course. Questions 1--4
are labelled with the objective they test; Question 5 spans everything.

\parttitle{Part 1: Module 4 Concept Review}

{\small

\noindent The whole module on one page. Everything in Part 2 can be answered from it.

\begin{enumerate}[label=\textbf{\arabic*.}, wide=0pt, itemsep=.4em, topsep=.5em]

\item \textbf{M4.L01 Bayesian inference.} The parameter gets a distribution: a \textbf{prior} before
the data, a \textbf{posterior} after, with
$f(\theta \mid x) \propto f(\theta)\,L(\theta;x)$. Feed each posterior in as the next prior to update
sequentially; the final answer does not depend on the order the data arrived in. When $P(H)$ is
small, one positive test can still leave $P(H \mid E)$ small, because false positives from the large
$H^c$ group swamp the true ones.

\item \textbf{M4.L02 Bayesian and frequentist.} A frequentist interval is a statement about the
\emph{procedure} across repeated samples; a \textbf{credible interval} is a statement about
$\theta$ given this prior and these data. A Bayesian test reports the posterior probability of a
claim, so there is no null distribution and no $p$-value.

\item \textbf{M4.L03 Conjugate models.} Beta--Binomial, Normal--Normal, Gamma--Poisson: the posterior
stays in the prior's family, so you update parameters instead of normalizing. In all three the
posterior mean is a weighted average,
$(w\cdot\text{prior mean} + n\cdot\text{data estimate})/(w+n)$, with prior weight $w$ equal to
$a+b$, $\sigma^2/\tau^2$, or $r$. The posterior always lands between the prior and the data.

\item \textbf{M4.L05 Approximation.} When no conjugate rule applies, use a \textbf{grid} (evaluate
prior $\times$ likelihood, normalize) or \textbf{MCMC}. A grid needs (points per axis)$^{k}$ cells, so
it collapses past a few parameters. For a chain: discard \textbf{burn-in}; a high acceptance rate is
\emph{not} a good sign; $\widehat{R}$ near $1.00$ says the chains agree; ESS is how many independent
draws you effectively have; $\text{MCSE} = \text{posterior sd}/\sqrt{\text{ESS}}$.

\item \textbf{M4.L04 With samples, everything is counting.} From posterior draws: the point estimate
is their mean or median, a $95\%$ credible interval is their $2.5$th and $97.5$th percentiles, and
$P(\theta > c)$ is the \emph{proportion} of draws above $c$. To predict a new observation, draw one
$y$ at \emph{each} posterior draw of $\theta$; that \textbf{posterior predictive} is always wider
than plugging in a single $\hat\theta$.

\item \textbf{M4.L06 Bayesian linear regression.} Same line, now with a prior on each of
$\beta_0$, $\beta_1$, and $\sigma$. \textbf{Check the prior before the data} with a \emph{prior
predictive check}: simulate parameters from the prior alone and look at the responses they imply. If
they are impossible, the prior is wrong no matter how ``uninformative'' it looks. After fitting, the
posterior is a compromise between prior and likelihood, and with a genuinely diffuse prior the
credible interval nearly coincides with the frequentist interval, even though the two mean different
things. Predicting still asks two questions: an interval for the \emph{mean line} at $x^{\ast}$, or a
wider one for \emph{one new case}.

\item \textbf{Four traps.} Calling a prior uninformative without checking what it implies. Reading a
posterior probability as a $p$-value. Predicting with a single $\hat\theta$ when the question was
about one new case. Trusting a chain because its acceptance rate is high.

\end{enumerate}

}

\newpage

\parttitle{Part 2: Question Set}

\noindent\textbf{Scenario.} For $n = 40$ penguins you model body \texttt{mass} in grams on
\texttt{flipper} length in mm. Flipper lengths run $172$ to $231$; masses run $2849$ to $6313$. The
likelihood alone (ordinary least squares) gives slope $47.640$ with standard error $3.335$,
$\hat\sigma = 380.6$, and a $95\%$ interval of $[\,40.889,\; 54.392\,]$.

\begin{questions}

%%%%%%%%%%%%%%%%%%%%%%%%%%%%%%%%%%%%%%
\item \textbf{Check the prior before the data.} \lo{M4.L06} Two candidate priors for the slope, each
paired with $\beta_0 \sim$ N$(0, 2000^2)$. Simulating from the prior \emph{alone} and reading off the
body mass implied at a flipper length of $200$ mm gives:

\begin{center}
\renewcommand{\arraystretch}{1.2}
\begin{tabular}{lc}
\toprule
\textbf{Prior on $\beta_1$} & \textbf{Implied mass at 200 mm, middle 95\%} \\
\midrule
N$(0, 100^2)$, called ``uninformative'' & $[\,-38{,}912,\;\; 41{,}024\,]$ g \\
N$(20, 5^2)$ & $[\,-337,\;\; 8{,}332\,]$ g \\
\bottomrule
\end{tabular}
\end{center}

\begin{parts}
\item Name the three parameters a Bayesian version of this model needs priors for.
\answerspace{3em}

\item The first prior is the one people reach for when they want to ``let the data speak.'' Point at
the single feature of its implied range that shows it is indefensible.
\answerspace{3.5em}

\item Both rows contain impossible values. Which row is nonetheless the more usable prior, and why?
\answerspace{4em}

\item Describe the prior predictive check as a recipe, and say why it must happen \emph{before} the
data are used.
\answerspace{4em}
\end{parts}

%%%%%%%%%%%%%%%%%%%%%%%%%%%%%%%%%%%%%%
\item \textbf{The posterior is a compromise.} \lo{M4.L06, M4.L03} Fitting with each prior gives:

\begin{center}
\renewcommand{\arraystretch}{1.25}
\begin{tabular}{lcccc}
\toprule
\textbf{Prior on $\beta_1$} & \textbf{Prior mean} & \textbf{Posterior mean} & \textbf{Posterior sd} & \textbf{95\% credible interval} \\
\midrule
N$(0, 100^2)$ & 0 & 47.587 & 3.333 & $[\,41.05,\; 54.12\,]$ \\
N$(20, 5^2)$ & 20 & 39.130 & 2.774 & $[\,33.69,\; 44.57\,]$ \\
\bottomrule
\end{tabular}
\end{center}

\noindent Under the first prior the prior carries weight $0.0011$ of the total; under the second it
carries $0.308$.

\begin{parts}
\item The first posterior mean is $47.587$ and the likelihood alone gave $47.640$. Explain why they
are so close, using the weight.
\answerspace{3.5em}

\item The second posterior mean, $39.130$, sits between $20$ and $47.640$. Verify that it is roughly
the weighted average the review describes, and say which way the prior pulled it.
\answerspace{4em}

\item The second posterior has the \emph{smaller} standard deviation. Is a narrower posterior
automatically better here? Explain.
\answerspace{4em}

\item A bootstrap of the same data gives a $95\%$ interval of $[\,42.06,\; 54.30\,]$, almost the
first row's credible interval. State what each of the two intervals claims, and why nearly equal
endpoints do not make them the same statement.
\answerspace{5em}
\end{parts}

%%%%%%%%%%%%%%%%%%%%%%%%%%%%%%%%%%%%%%
\item \textbf{Two questions at one flipper length.} \lo{M4.L04} At $\texttt{flipper} = 200$ the
fitted mass is $4753.8$ g, with $\text{SE}_{\text{fit}} = 61.7$ and $\hat\sigma = 380.6$.

\begin{parts}
\item One interval is $[\,4628.8,\; 4878.7\,]$ and the other is $[\,3973.3,\; 5534.3\,]$. Label each
with the question it answers.
\answerspace{3.5em}

\item Compute $\sqrt{\hat\sigma^2 + \text{SE}_{\text{fit}}^2}$ and use it to say which term makes
the second interval about six times as wide.
\answerspace{4em}

\item In the Bayesian version, both intervals come from the same posterior draws. Describe the two
different recipes, in terms of what you do at each draw.
\answerspace{4.5em}

\item A colleague catches one penguin with a $200$ mm flipper and reports the narrower interval.
What have they claimed, and what should they have reported?
\answerspace{3.5em}
\end{parts}

%%%%%%%%%%%%%%%%%%%%%%%%%%%%%%%%%%%%%%
\item \textbf{The rest of Module 4.} \lo{M4.L01, L03, L05} Short answers.

\begin{parts}
\item A detector with sensitivity $0.92$ and specificity $0.88$ fires at a site where the species is
present $3\%$ of the time. Without computing, say whether the posterior probability is nearer $0.2$
or $0.9$, and name the quantity that decides it.
\answerspace{3.5em}

\item A Beta$(2,8)$ prior meets $5$ successes in $20$ trials. Give the posterior and its mean, and
say what the prior was worth in trials.
\answerspace{3.5em}

\item Two samplers on the same posterior: one accepts $98.7\%$ of proposals with ESS $14$, the other
accepts $61\%$ with ESS $867$. Which do you keep, and what is wrong with the other?
\answerspace{4em}

\item Four chains give $\widehat{R} = 1.57$. What does that number compare, and what must you do
before reporting anything?
\answerspace{3.5em}

\item With posterior sd $0.586$ and ESS $867$, compute the MCSE. Then say what ESS you would need to
halve it.
\answerspace{3.5em}
\end{parts}

%%%%%%%%%%%%%%%%%%%%%%%%%%%%%%%%%%%%%%
\item \textbf{The whole course, side by side.} \lo{all four modules} Complete the translation table.

\begin{center}
\renewcommand{\arraystretch}{1.9}
\begin{tabularx}{\linewidth}{LLL}
\toprule
\textbf{Question} & \textbf{Frequentist answer} & \textbf{Bayesian answer} \\
\midrule
What is treated as random? & & \\
A range for a parameter & & \\
Evidence against a claim & & \\
Where prior knowledge enters & & \\
Predicting one new case & & \\
\bottomrule
\end{tabularx}
\end{center}

\begin{parts}
\item Two rows have answers that can be numerically almost identical while meaning different things.
Which two, and what is the difference in what is being claimed?
\answerspace{4.5em}

\item A classmate's final report ends with the four sentences below. Mark each \textbf{defensible} or
\textbf{not}, and rewrite the ones that are not.
\begin{enumerate}[label=\textbf{S\arabic*.}, leftmargin=2.2em, itemsep=.25em, topsep=.35em]
\item ``Our $95\%$ confidence interval for the slope is $[40.9, 54.4]$, so there is a $95\%$
probability the true slope is in that range.''
\item ``We used an uninformative N$(0, 100^2)$ prior, so our conclusions are not affected by any
prior beliefs.''
\item ``The posterior probability that $\beta_1 > 45$ is $0.78$, which is not significant at the
$5\%$ level.''
\item ``Our credible interval and our bootstrap interval nearly agree, which proves the prior was
correct.''
\end{enumerate}
\answerspace{7em}
\end{parts}

\end{questions}

\end{document}
